% Part: proof-theory
% Chapter: sequent-calculus
% Section: introduction

\documentclass[../../../include/open-logic-section]{subfiles}

\begin{document}

\olfileid[id]{pt}{seq}{int}

\olsection{Pendahuluan}

Kalkulus sekuen adalah sistem !!{proof} yang diperkenalkan pada dekade 1930-an
oleh Gerhard Gentzen. Sistem ini tidak terutama dimaksudkan sebagai metode
praktis untuk !!{proving} sesuatu. Sebaliknya, Gentzen menggunakannya untuk
membuktikan hasil-hasil tertentu tentang kemungkinan struktur !!{proof}s dan
sifat-sifatnya. Pertanyaan semacam inilah yang menjadi pokok perhatian teori
bukti.

Sebagaimana tersirat dari namanya, kalkulus sekuen bekerja dengan
\emph{sekuen}, bukan dengan !!{formula}s. Sekuen adalah pasangan barisan atau
multihimpunan !!{formula}s. Sistem asal Gentzen (\Log{LK} dan \Log{LJ})
menggunakan barisan; para ahli teori bukti kini lebih menyukai multihimpunan.
Multihimpunan menyerupai himpunan, tetapi suatu !!{element} dapat muncul lebih
dari sekali di dalam multihimpunan. Namun, seperti dalam himpunan, urutan
!!{element}s tidak penting. Jadi, ekspresi $\{!A,!A,!B\}$ dan
$\{!A,!B,!A\}$ mendeskripsikan multihimpunan yang sama, tetapi multihimpunan
itu berbeda dari $\{!A,!B\}$ dan $\{!A,!A,!B,!B\}$, karena dua yang terakhir
memuat jumlah $!A$ dan $!B$ yang berbeda dari yang pertama.

Secara tradisional, multihimpunan !!{formula}s dalam teori bukti ditulis
dengan huruf Yunani kapital. Jadi, ketika kita menulis $\Gamma$, $\Delta$,
$\Pi$, $\Lambda$, atau $\Theta$, kita memaksudkannya sebagai multihimpunan
!!{formula}s dari logika yang sedang digunakan. Secara tradisional pula,
alih-alih menulis $\tuple{\Gamma,\Delta}$, para ahli teori bukti memisahkan
kedua komponennya dengan simbol sekuen; kita menggunakan~$\Sequent$.

\begin{defn}[Sekuen]
Suatu \emph{sekuen} adalah ekspresi berbentuk
\[
\Gamma \Sequent \Delta
\]
dengan $\Gamma$ dan $\Delta$ sebagai multihimpunan berhingga (mungkin kosong)
dari !!{formula}s. $\Gamma$ disebut \emph{anteseden}, sedangkan $\Delta$
disebut \emph{sukseden}.
\end{defn}

Misalkan $!G$ adalah konjungsi !!{formula}s dalam~$\Gamma$ (dengan konjungsi
kosong dianggap sebagai~$\ltrue$), dan $!D$ adalah disjungsi !!{formula}s
dalam~$\Delta$ (dengan disjungsi kosong dianggap sebagai~$\lfalse$). Maka
sekuen $\Gamma\Sequent\Delta$ benar (dalam suatu struktur~$\Struct{M}$) jika
dan hanya jika $!G\lif !D$ benar. Dalam logika klasik, hal ini sama dengan
mengatakan bahwa suatu !!{formula} dalam~$\Gamma$ salah atau suatu
!!{formula} dalam~$\Delta$ benar.

Jika $\Gamma$ merupakan multihimpunan !!{formula}s, kita menulis
$\Gamma,!A$ (atau $!A,\Gamma$) bagi hasil penambahan $!A$ ke~$\Gamma$.
Jika $\Delta$ juga merupakan multihimpunan !!{formula}s, maka
$\Gamma,\Delta$ adalah gabungan multihimpunan keduanya---jika $\Gamma$
memuat $!A$ dengan \emph{multiplisitas}~$n$ (yakni, memuat $n$ salinan~$!A$)
dan $\Delta$ memuat $!A$ dengan multiplisitas~$m$, maka $\Gamma,\Delta$
memuat $!A$ dengan multiplisitas~$n+m$. Jika multihimpunan pada salah satu
sisi sekuen kosong, biasanya kita membiarkan sisi itu kosong. Misalnya,
$\Sequent !A$ adalah sekuen dengan anteseden kosong dan sukseden yang memuat
$!A$ dengan multiplisitas~$1$.

!!^a{proof} dalam kalkulus sekuen adalah pohon sekuen, dengan sekuen-sekuen
teratas (atau awal) sebagai aksioma dalam sistem yang ditinjau; jika suatu
sekuen terletak di bawah satu atau dua sekuen lain, sekuen tersebut harus
mengikuti secara benar melalui suatu aturan inferensi dalam sistem. Mula-mula
kita akan meninjau dua sistem sekuen berbeda bagi logika klasik,
$\Log{G1c}$ dan~$\Log{G3c}$. Aksioma dan aturannya diberikan dalam
\cref{pt:seq:int:tab:G1c,pt:seq:int:tab:G3c}.
\oliflabeldef{fol:seq::chap}{ (Sistem asal Gentzen \Log{LK} dideskripsikan
dalam \olref[fol][seq][]{chap}.)}{}

\subfile{rules-G1c}
\subfile{rules-G3c}

Kedua sistem berbeda secara halus, dan perbedaan tersebut menyebabkan sifat
teori-bukti yang berbeda. \Log{G1c} mempunyai aksioma $!A\Sequent !A$ tanpa
!!{formula} lain yang diizinkan. \Log{G3c} mempunyai aksioma berbentuk
$!C,\Gamma\Sequent\Delta,!C$ dengan !!{formula}s ``samping'' sebarang dalam
$\Gamma$ dan~$\Delta$, tetapi !!{formula}~$!C$ yang muncul pada kedua sisi
harus \emph{atomik}. \Log{G1c} masing-masing mempunyai dua versi aturan
$\LeftR{\land}$ dan $\RightR{\lor}$, sedangkan \Log{G3c} masing-masing hanya
mempunyai satu. Selain itu, \Log{G1c} mempunyai ``aturan struktural''
tambahan yang disebut ``kontraksi''
(\LeftR{\Contraction}, \RightR{\Contraction}) dan ``pelemahan''
(\LeftR{\Weakening}, \RightR{\Weakening}).

Sistem \Log{G1c} dan \Log{G3c} sama-sama sahih dan lengkap bagi logika
klasik. Dengan kata lain, setiap sekuen yang !!{provable} dalam sistem ini
benar dalam setiap !!{structure}, dan setiap sekuen yang benar dalam setiap
!!{structure} mempunyai !!a{proof}. Jadi, pertanyaan \emph{apakah} suatu
sekuen mempunyai !!a{proof} juga dapat dijawab dengan metode logika lain
(misalnya metode teori model). Karena kedua sistem !!{prove} tepat sekuen
yang benar dalam setiap !!{structure}, keduanya secara khusus membuktikan
sekuen-sekuen yang sama.

Namun, teori bukti tidak hanya membahas \emph{apakah} sesuatu dapat
dibuktikan, tetapi juga \emph{bagaimana}. Para ahli teori bukti mengajukan
pertanyaan seperti ``Dapatkah aturan tertentu selalu dihindari?'', ``Dapatkah
aturan ini ditambahkan pada sistem tanpa membuat sekuen baru menjadi dapat
dibuktikan?'', atau ``Dapatkah !!{proof}s dalam satu sistem diubah menjadi
!!{proof}s dalam sistem lain?'' Jika jawaban atas pertanyaan semacam itu
adalah ``ya'', faktanya sering dibuktikan melalui induksi pada kompleksitas
!!{proof}s, atau secara ekuivalen melalui induksi pada struktur !!{proof}s.
Hal ini tidak hanya menghasilkan bukti bagi fakta tersebut (misalnya, jika
\Log{G1c} membuktikan suatu sekuen, maka \Log{G3c} membuktikan sekuen yang
sama), tetapi juga suatu prosedur (misalnya, untuk mengubah !!{proof}s dalam
\Log{G1c} menjadi !!{proof}s dalam~\Log{G3c}).

Salah satu hasil utama teori bukti adalah \emph{teorema eliminasi potongan}.
Teorema itu menunjukkan bahwa kita dapat menambahkan aturan potongan
\[
\Axiom$\Gamma \fCenter \Delta, !A$
\Axiom$!A, \Pi \fCenter \Lambda$
\RightLabel{\Cut}
\BinaryInf$\Gamma, \Pi \fCenter \Delta, \Lambda$
\DisplayProof
\]
pada sistem sekuen tanpa mengubah sekuen mana yang !!{provable}. Lebih jauh,
buktinya memberikan metode yang mengubah setiap !!{proof} yang menggunakan
aturan \Log{G1c} (atau \Log{G3c}) bersama \Cut{} menjadi bukti yang hanya
menggunakan aturan \Log{G1c} (atau \Log{G3c}). Dengan kata lain, prosedur
ini ``mengeliminasi'' aturan \Cut{} dari !!{proof}s yang menggunakannya;
karena itulah namanya.

\end{document}
